\section{Conclusiones y recomendación}

\textbf{Qué mejoró.} Bajo los supuestos del modelo, redistribuir las mismas 23 unidades de
capacidad con el horario candidato (1, 2, 5, 5, 6, 4) reduce la espera hasta el inicio de la
atención respecto del horario base (3, 3, 4, 4, 5, 4).

\textbf{Cuánto.} Con demanda nominal, la espera media simulada baja de \WaitBaseNom{} a
\WaitCandNom{} min (reducción pareada de \RedNom{} min, \RedPctNom\,\%; IC 95\,\%
[\RedCILowNom, \RedCIHighNom] min), y también bajan la mediana, el percentil 90, la
proporción con espera mayor a 60 minutos y la cola.

\textbf{Bajo qué condiciones.} El efecto depende de la carga: con demanda alta la
reducción absoluta es mayor (\RedHigh{} min), y con demanda baja no se distingue de cero
(IC 95\,\% [\RedCILowLow, \RedCIHighLow] min). La mejora es consistente solo cuando el
sistema está lo bastante exigido.

\textbf{Qué tan confiable es la comparación.} La comparación entre horarios es pareada, con
semillas independientes de las de optimización y \NReps{} réplicas por condición, pero la
variabilidad entre días es alta. La validación frente a NHAMCS es parcial: el modelo base
sobreestima la espera observada (MAE de \MaeBase{} min en la base y de \MaeCand{} min en la
candidata), por lo que las cifras absolutas de espera no deben tomarse como predicciones.
Con más confianza se puede leer la \emph{dirección} del efecto que su magnitud.

\textbf{Qué limitaciones permanecen.} El servicio es un supuesto; no hay comparador de
optimización ni sensibilidad a los parámetros del GA y del servicio; el modelo es un recurso
agregado calibrado con un DE nacional estadounidense, no con un hospital guatemalteco.

\textbf{Recomendación.} Conservar la candidata como escenario de análisis y demostración,
no como política lista para aplicar. Como pasos siguientes, en orden de prioridad:
(1)~calibrar el servicio con datos de ocupación del equipo (o al menos propagar la
incertidumbre del supuesto con un análisis de sensibilidad); (2)~añadir una búsqueda
aleatoria con el mismo presupuesto para medir el aporte del GA; (3)~ampliar la
sensibilidad a demanda, presupuesto y pesos del objetivo; (4)~validar con datos locales
antes de cualquier decisión operativa.

\begin{alcance}
La conclusión es sobre una \textbf{mejora simulada} bajo supuestos explícitos y una
referencia pública de Estados Unidos. No demuestra causalidad ni una mejora clínica real.
\end{alcance}
